\chapter{Cơ sở lý thuyết}
\label{chap:co-so}

\section{Định vị lỗi dựa trên phổ độ phủ}
\label{sec:sbfl}

\subsection{Ma trận độ phủ và các công thức nghi ngờ}

Cho chương trình $P$ và bộ kiểm thử $T = T_f \cup T_p$ gồm các kiểm thử thất bại $T_f$ và thành công $T_p$. Với mỗi
dòng $s$, ký hiệu $e_f(s)$, $e_p(s)$ là số kiểm thử thất bại, thành công chạy qua $s$; $n_f(s)$, $n_p(s)$ là số kiểm
thử thất bại, thành công không chạy qua $s$. Ma trận độ phủ ghi lại, với mỗi cặp (kiểm thử, dòng), kiểm thử đó có chạy
qua dòng đó hay không. SBFL gán cho mỗi dòng một điểm nghi ngờ; ba công thức dùng trong tiểu luận là
Tarantula~\cite{jones2002}, Ochiai~\cite{abreu2009} và DStar~\cite{wong2014dstar}:
\begin{align}
\mathrm{Tarantula}(s) &= \frac{e_f/(e_f+n_f)}{e_f/(e_f+n_f) + e_p/(e_p+n_p)}, \label{eq:tarantula}\\
\mathrm{Ochiai}(s) &= \frac{e_f}{\sqrt{(e_f+n_f)\,(e_f+e_p)}}, \label{eq:ochiai}\\
\mathrm{DStar}(s) &= \frac{e_f^{\,*}}{e_p + n_f}, \quad * = 2. \label{eq:dstar}
\end{align}
Các dòng được xếp giảm dần theo điểm; dòng không được kiểm thử thất bại nào chạy qua ($e_f = 0$) có điểm 0 và trong
tiểu luận không được xếp hạng.

\subsection{Chỉ số đánh giá và cách xử lý hòa điểm}

Với mỗi chương trình, \emph{hạng} của lỗi là hạng tốt nhất của một dòng lỗi chuẩn trong danh sách xếp hạng.
Acc@$k$ là tỉ lệ chương trình có hạng không vượt quá $k$:
\begin{equation}
\mathrm{Acc@}k = \frac{1}{|\mathcal{P}|} \sum_{p \in \mathcal{P}} \mathbb{1}\!\left[\,\mathrm{rank}_p \le k\,\right],
\label{eq:acck}
\end{equation}
trong đó $\mathcal{P}$ là tập chương trình \emph{cố định}: chương trình mà công cụ không chạy được vẫn nằm trong mẫu số
và được tính là trượt. Khi nhiều dòng cùng điểm, có ba cách phá hòa: lấy vị trí đầu (lạc quan, theo thứ tự dòng), vị
trí cuối (bi quan) hoặc \emph{hạng trung bình} của nhóm hòa. Tiểu luận dùng hạng trung bình, theo khuyến nghị của các
nghiên cứu về đánh giá định vị lỗi~\cite{pearson2017,desouza2016}; cách chấm của tác giả SemLoc dùng vị trí đầu.

\section{Lỗi ngữ nghĩa}
\label{sec:loi-ngu-nghia}

\subsection{Vì sao độ phủ mất tín hiệu}

Lỗi ngữ nghĩa là lỗi làm chương trình tính sai giá trị trong khi vẫn chạy qua cùng các câu lệnh như khi đúng --- ví
dụ \texttt{updated = cap - 1} thay vì \texttt{updated = cap}. Khi cả kiểm thử thành công và thất bại đều chạy qua dòng
lỗi, $e_p(s) > 0$ làm điểm Ochiai của dòng lỗi giảm và ngang với các dòng xung quanh; SBFL khi đó gần như đoán mò trong
phạm vi hàm.

\subsection{Định nghĩa khách quan bằng độ phủ}

Để phân tích kết quả theo loại lỗi, tiểu luận dùng một định nghĩa tự động, chỉ dựa trên độ phủ của các dòng lỗi chuẩn:
\begin{itemize}
\item \textbf{lỗi ngữ nghĩa} (semantic): có dòng lỗi được ít nhất một kiểm thử thất bại \emph{và} một kiểm thử thành
công chạy qua --- độ phủ không phân biệt được;
\item \textbf{lộ qua độ phủ} (coverage-revealed): kiểm thử thất bại chạy tới dòng lỗi nhưng không kiểm thử thành công
nào chạy tới;
\item \textbf{không thực thi} (not executed): không kiểm thử thất bại nào chạy tới dòng lỗi chuẩn --- thường do dòng
lỗi chuẩn bị đặt sai hoặc là dòng trống.
\end{itemize}
Nhãn được tính lại theo đúng bộ dòng lỗi chuẩn đang dùng để chấm (gốc hoặc đã kiểm chứng).

\section{SemLoc: phổ vi phạm ràng buộc ngữ nghĩa}
\label{sec:semloc}

\subsection{Quy trình bốn bước}

Hình~\ref{fig:semloc} tóm tắt quy trình của SemLoc~\cite{yang2026semloc}. Mỗi ràng buộc $c$ do LLM sinh ra là một bộ
$\langle$mã, loại, vùng, điểm neo, biểu thức, ý đồ$\rangle$; biểu thức được chèn vào chương trình dưới dạng một lời gọi
kiểm tra tại điểm neo. Chạy bộ kiểm thử cho ma trận ``ràng buộc $\times$ kiểm thử'' ghi ràng buộc nào bị vi phạm trong
kiểm thử nào; thay ``dòng'' bằng ``ràng buộc'' trong công thức~\eqref{eq:ochiai} cho điểm của từng ràng buộc, và điểm
của một dòng là điểm lớn nhất trong các ràng buộc được quy về dòng đó. Bước kiểm chứng phản thực tế loại ràng buộc bị
vi phạm cả trong kiểm thử thành công, rồi cho LLM vá thử tại dòng được quy cho từng ràng buộc còn lại và chạy lại kiểm thử.

\begin{figure}[H]
\centering
{\setstretch{1}%
\begin{tikzpicture}[node distance=4mm and 6mm, font=\small,
  box/.style={draw, rounded corners, align=center, minimum height=11mm, text width=24mm, fill=blue!4},
  ours/.style={draw=orange!80!black, dashed, rounded corners, align=center, minimum height=11mm, text width=30mm,
               font=\footnotesize, fill=orange!6},
  arr/.style={-{Stealth[length=2mm]}}]
\node[box] (a) {Chương trình\\+ kiểm thử};
\node[box, right=of a] (b) {LLM sinh\\ràng buộc};
\node[box, right=of b] (c) {Chèn mã,\\chạy kiểm thử};
\node[box, right=of c] (d) {Phổ vi phạm\\(Ochiai)};
\node[box, below=10mm of d] (e) {Kiểm chứng\\phản thực tế};
\node[box, left=of e] (f) {Xếp hạng\\dòng};
\node[ours, left=of f] (g) {+\ConstPriorWeight{} điểm cho dòng LLM\\tự đoán (không có\\trong bài báo)};
\draw[arr] (a) -- (b); \draw[arr] (b) -- (c); \draw[arr] (c) -- (d); \draw[arr] (d) -- (e); \draw[arr] (e) -- (f);
\draw[arr, dashed, orange!80!black] (b) -- (b |- g.north); \draw[arr, dashed, orange!80!black] (g) -- (f);
\end{tikzpicture}}
\caption{Quy trình của SemLoc (theo~\cite{yang2026semloc}); khung nét đứt là thành phần nhóm phát hiện trong mã nguồn}
\label{fig:semloc}
\end{figure}

\subsection{Điểm thưởng cho dòng LLM tự đoán}

Khi đọc mã nguồn công khai, nhóm thấy một thành phần không được mô tả trong bài báo: ngoài các ràng buộc, LLM còn trả
về một dòng lỗi mà nó tự đoán (trường \texttt{fault\_line}), và dòng này được cộng thêm \ConstPriorWeight{} điểm --- lớn hơn điểm tối đa
1 của một ràng buộc. Vì vậy, dòng LLM tự đoán \emph{luôn} đứng hạng 1, bất kể phổ vi phạm ra sao. Thành
phần này là trọng tâm của câu hỏi ``grounded hay guessed'' ở các chương sau.

\section{Phân tích tĩnh và phân tích động}
\label{sec:tinh-dong}

Phân tích tĩnh xét mã nguồn mà không chạy chương trình. Trình kiểm tra phong cách mã (lint) như ruff phát hiện mẫu mã
dễ lỗi hoặc khó đọc (biến không dùng, so sánh thừa, cách viết dài dòng); trình kiểm tra kiểu như pyrefly suy kiểu và báo
chỗ kiểu không khớp với chú thích. Cả hai chỉ biết mã \emph{được viết} thế nào, không biết mã \emph{được mong đợi} làm
gì: viết \texttt{>} thay cho \texttt{>=} hay trừ phí hai lần đều hợp lệ về cú pháp và về kiểu. Phân tích động lấy thông
tin từ việc chạy chương trình: kết quả kiểm thử, độ phủ, giá trị biến. Với lỗi ngữ nghĩa, chỉ kiểm thử --- bản đặc tả do
lập trình viên viết --- mới chứa thông tin về ý định, nên định vị lỗi ngữ nghĩa thuộc về phân tích động; phân tích tĩnh
\IfAtLeastElse{lan3}{được dùng như một mốc so sánh (Chương~\ref{chap:ket-qua})}{sẽ được dùng như một mốc so sánh ở giai đoạn sau}.

\section{Công cụ thống kê}
\label{sec:thong-ke}

\subsection{Kiểm định cặp}

Mọi so sánh giữa hai phương pháp đều là so sánh \emph{cặp} trên cùng các chương trình. Với chỉ số nhị phân (trúng
hay trượt top-$k$), kiểm định McNemar dùng số chương trình chỉ phương pháp A trúng ($b$) và chỉ phương pháp B trúng ($c$),
bản chính xác khi $b+c$ nhỏ. Với hạng nghịch đảo ($1/\mathrm{rank}$, trượt là 0), kiểm định dấu hạng Wilcoxon cùng
hệ số tương quan hạng hai chuỗi ghép cặp làm độ lớn hiệu ứng:
\begin{equation}
r_{rb} = \frac{R^{+} - R^{-}}{R^{+} + R^{-}},
\label{eq:rankbiserial}
\end{equation}
với $R^{+}$, $R^{-}$ là tổng hạng của các chênh lệch dương, âm (bỏ chênh lệch bằng 0).

\subsection{Khoảng tin cậy bootstrap theo cụm và giá trị Shapley}

Mỗi cấu hình được chạy ba lần. Kết quả của một chương trình được lấy trung bình qua ba lần \emph{trước khi} lấy mẫu
lại, nên đơn vị lấy mẫu là chương trình chứ không phải lần chạy. Vì nhiều chương trình của SemFault-250 sinh ra từ cùng
một khuôn mẫu, khoảng tin cậy \ConstCiLevel\% được tính bằng bootstrap \emph{theo cụm khuôn mẫu} (lấy lại cả cụm), với \ConstBootstrapDraws{} lần
lấy mẫu và hạt giống cố định.

Để đo phần đóng góp của từng tín hiệu, tiểu luận dùng giá trị Shapley. Với tập tín hiệu $N$ và hàm giá trị $v(S)$ là
Acc@$k$ khi chỉ bật các tín hiệu trong $S$,
\begin{equation}
\phi_i = \sum_{S \subseteq N\setminus\{i\}} \frac{|S|!\,(|N|-|S|-1)!}{|N|!}\,\bigl(v(S\cup\{i\}) - v(S)\bigr),
\label{eq:shapley}
\end{equation}
và $\sum_i \phi_i = v(N)$. Vì Acc@$k$ là trung bình các giá trị trúng/trượt theo chương trình, giá trị Shapley tổng bằng
trung bình giá trị Shapley của từng chương trình; tính chất này cho phép tính khoảng tin cậy bootstrap cho từng $\phi_i$.
